\section{Conclusion}
\label{sec:conclusion}

Deployed anomaly detectors are thresholds, not scores, and a threshold calibrated on one domain breaks on the next. Across six leave-one-domain-out folds over four domain families of industrial time series, static source-calibrated thresholds failed in both directions---over-alarming to $0.18$ coverage on wind data and under-alarming to $0.999$ elsewhere---while an online adaptive-conformal layer restored normal-window coverage to within $0.015$ of nominal on five folds at $0.6$--$0.8\,\mu$s per decision. Target-aware refinements (divergence-shifted, conditional) brought no gain, divergence did not predict the $0.5\times$--$19\times$ transfer cost, and on two folds whose scores rank anomalies below chance, calibration controlled coverage but could not make the alarms useful. We release the benchmark, protocols, and code, and we state the negative results plainly: deploy calibration, audit the ranker first, and never trust a source threshold on a new domain.
